\begin{titlepage}
    \begin{abstract}
        La gestione tradizionale di contest, basata su strumenti eterogenei, è un processo frammentato e insicuro che compromette l’integrità dei dati e il coordinamento in tempo reale.
        L'analisi dello stato dell'arte ha confermato un \textbf{vuoto di mercato} per soluzioni verticali, motivando la progettazione di \textbf{Swift Contest}: una piattaforma software \emph{end-to-end}, sicura e multipiattaforma.

        Per soddisfare i requisiti emersi dall’analisi, è stato selezionato uno stack tecnologico moderno.
        L’esigenza di sviluppo \emph{cross-platform} ha guidato la scelta del framework \textbf{Flutter}, mentre la necessità di un backend robusto e relazionale ha portato all’adozione di \textbf{Supabase (PostgreSQL)}.
        La logica serverless è stata implementata in modo strategico: le \textbf{Funzioni RPC} in PL/pgSQL garantiscono l' \textbf{atomicità delle transazioni}, mentre le \textbf{Edge Functions} in TypeScript gestiscono le \textbf{interazioni con API esterne e con lo Storage}.

        La progettazione del sistema si fonda su un’architettura sicura basata sul \textbf{Principio del Minimo Privilegio (PoLP)}, dove la \textbf{Row Level Security (RLS) e un'API controllata} lavorano in sinergia per mediare ogni accesso ai dati.
        Sul frontend, il pattern \textbf{MVVM con BLoC} garantisce la \textbf{Separazione delle Responsabilità (SoC)}, mentre sul backend, il \textbf{meccanismo di snapshot} assicura l’integrità del voto.
        Un sistema di \textbf{autenticazione ibrido} supporta inoltre l’accesso sicuro per ``Giurati Ospiti'' tramite autenticazione anonima.

        La fase finale del progetto ha visto la distribuzione delle sue componenti su piattaforme dedicate: le migrazioni IaC e le funzioni serverless sull'infrastruttura \textbf{Supabase}, il client web su \textbf{Vercel} tramite pipeline di CI/CD e l'artefatto Android su \textbf{GitHub}.
        Questo processo ha portato alla realizzazione di un \textbf{prototipo completo e funzionante}.
        In conclusione, Swift Contest non solo valida la fattibilità dell'approccio architetturale, ma si configura come un \textbf{blueprint robusto e scalabile} per la creazione di applicazioni collaborative sicure, dimostrando come un ecosistema tecnologico moderno possa colmare una reale esigenza di mercato.
    \end{abstract}
\end{titlepage}